% !TEX root = ../oddp.tex

\section{Additional remarks}

\subsection{Homotopy retractions}\label{s:hretractions} In the previous section, one could try to define a homotopy retraction of $\Waug{r}$, and then extend the map between the augmentations of $\Per{r}$ and $\Waug{r}$ (which is the identity map $R\to R$) using this homotopy retraction. This would solve our problem if in addition the homotopy retraction is periodic with respect to the endomorphism $\theta$ of $\Waug{r}$. For example, if $r=3$, one can use the following retraction $\eta$:
\begin{align*}
	\eta(1) &= \rho^0e_1 
	\\
	\eta(\rho^0e_{2k-1}) &= 0
	&
	\eta(\rho^1e_{2k-1}) &= \rho^0e_2
	&
	\eta(\rho^2e_{2k-1}) &= -\rho^2e_2
	\\
	\eta(\rho^0e_{2k}) &= 0 
	&
	\eta(\rho^1e_{2k}) &= (\rho^0+\rho^1+\rho^2)e_3
	&
	\eta(\rho^2e_{2k}) &= 0 
\end{align*}  
and define recursively 
\begin{align*}
	\phi(\sus{k}[0]) &= \eta\circ \phi \circ \partial(\sus{k}[0])
	\\
	\phi(\sus{k}[0,1]) &= \eta\circ \phi \circ \partial(\sus{k}[0,1])
\end{align*}
and extend equivariantly. Finding such a periodic homotopy retraction in general seems as challenging as finding the map $\phi$.


\subsection{The bar resolution}\label{ss:milgram} Recall the that the Milgram construction applied to a group $G$, yields a contractible augmented simplicial set $\EG$ with a free $G$-action:
\begin{align*}
	(\EG)_k &= G^{k+1}\\
	\d_i(g_0\dots g_k) &= (g_0\dots\hat{g}_i\dots g_k) \\
	\s_i(g_0\dots g_k) &= (g_0\dots g_i,g_i\dots g_k)
\end{align*}
We can regard $\EC_r$ as the union $\bigcup_n \chains(\Cyc_r)^{\ot n}$ (this filters $\EC_r$ by degree). Then, the inclusion of $\Cyc_r$ as the vertices of $\asimplex^{r-1}$ induces a map $\chains(\EC_r) \to \bigcup_n\chains(\asimplex^{r-1})^{\ot n}$ that sends a tuple $(a_0,a_1,\dots,a_k)$ to the generator $[a_0]\ot \ldots \ot [a_k]$. The composition
\[
\chains(\EC_r) \lra \bigcup_n\chains(\asimplex^{r-1})^{\ot n} \lra \Waug{r}
\]
is an equivariant chain map. We can describe the dual of this map as follows: A generator $(-1)^{s}\cdot c\cdot (a_0,\ldots,a_{q-1})$ is \emph{admissible} if 
\begin{itemize}
	\item The generator $(-1)^{s'}\cdot (c\cdot \tilde{r}!) [a_0,\ldots,a_{q-r}]$ is admissible and $[a_{q-r},\ldots,a_{k-1}]$ does not contain repetitions and $(-1)^{s'+s}$ computes the parity of this permutation.
	\item $q<r$, then $s$ is the sign of the permutation that orders $[a_0,\ldots,a_{q-1}]$, $c$ equals $\frac{\varphi(q)}{\tilde{r}!}$, the number $r-1$ does not belong to the word, and the word obtained after ordering it starts with an even entry and alternates even and odd entries.
\end{itemize}

%
 %This composition has a description analogous to the one used to define $\Psiom$, using the following suspension map:
%\[
%S(a_0,\dots,a_k) = \begin{cases}
	%(-1)^{\sign{\perm}} (a_{0},\dots,a_{k-r+1}) & \text{ if $k \geq r-1$} \\
	%\emptyset & \text{if $k = r-2$ and $a_i\neq a_j$ for all $i,j$} \\
	%0 & \text{otherwise.}
%\end{cases}
%\]
%where $\pi$ is the permutation that orders $(a_{k-r+1},\dots,a_{k})$. The image of $e_{-q}^{\vee}$ is the sum of all words $A = (a_0,\dots,a_{q-1})\in \cochains[-q] \EC_r$ such that
%\begin{itemize}
%\item none of the subwords of the form $(a_{q-1-\ell(r-1)},a_{q-1-(\ell-1)(r-1))}$ of length $r$ contains repeated elements,
%\item the only subword of the form $(a_0,\dots,a_{q-1-\ell(r-1)})$ of length smaller than $r$ becomes, after having been ordered, an increasing sequence that alternates even and odd entries different from $r-1$, starting with an even entry.
%\end{itemize}
%each word has the sign needed to order each of the subwords mentioned above, and the coefficient $\frac{\varphi(m)!}{(\tilde{r}!)^j}$ where $m$ is the length of the last of the subwords mentioned, and $j$ is the total number of subwords and $\varphi(m) = \lfloor\frac{r-m-1}{2}\rfloor$.
\begin{example} If $r=3$,
\begin{align*}
	\Psi^\vee(e_{-1}^\vee) &= (0)
	\\
	\Psi^\vee(e_{-2}^{\vee}) &= (0,1)-(1,0)
	\\
	\Psi^\vee(e_{-3}^{\vee}) &= (0,1,2) - (0,2,1)
\\
	\Psi^\vee(e_{-4}^{\vee}) &= (0,1,2,0) - (0,1,0,2) + (1,0,2,1) - (1,0,1,2)
\end{align*}
if $r= 5$,
\begin{align*}
	\Psi^\vee(e^\vee_{-1}) &= \frac{1}{2}\left((0) + (2)\right)
	\\
	\Psi^\vee(e^\vee_{-2}) &= \frac{1}{2}\left((0,1) - (1,0) + (0,3) - (3,0) + (2,3) + (3,2)\right)
	\\
	\Psi^\vee(e^\vee_{-3}) &= \frac{1}{2}\left((0,1,2) -(0,2,1) + (1,2,0) - (1,0,2) + (2,0,1) - (2,1,0)\right)
\end{align*}
$\Psi^\vee(e^\vee_{-4})$ is the signed sum of all permutations of $(0,1,2,3)$, so it has $24$ summands, again with coefficient $\frac{1}{2}$.
\end{example}

This map is different from the one obtained in \cite[Prop.~6.16]{brumfiel2023explicit}.



\subsection{The even prime}\label{ss:even} We have only constructed a map $\phi$ when $r$ is odd, but if $r= 2$ and $R=\bF_2$, then $\Per{r}$ and $\Waug{r}$ are isomorphic and we can define an isomorphism $\phi\colon \Per{r}\to \Waug{r}$ sending the generator $\sus{q-1}[0]$ to $e_q$.

Regarding the map $\psi$, observe that $\partial \simplex^{r-1}$ is isomorphic to $\Cyc_r$, thus the union $\bigcup_n \chains(\partial \simplex^1)^{\ot n}$ is isomorphic to the unnormalized chain complex of the Milgram resolution $\EC_2$, while $\Per{r}$ is isomorphic to the normalized chain complex of $\EC_2$. The map $\psi\colon \chains(\partial \asimplex^1)^{\ot (n+1)} \lra \Per{r}$ is the composition of
\[
	\chains(\partial \asimplex^1)^{\ot (n+1)}
	\to
	\uchains(\EC_r)
	\to
	\chains(\EC_r)
	\overset{\cong}{\to}
	\Per{r}
\] 
where the first map forgets the empty factors, the second map is the projection to the normalised chains and the last map sends a normalised chain $[u_0,u_1,\ldots,u_{q-1}]$ to $\sus{q-1}[u_0]$. 

Under this identification, the coface map $\d^i$ of $\chains(\partial\simplex^{r-1})^{n+1}$ (resp.\ $\chains(\partial\asimplex^{r-1})^{n+1}$) sends a sequence $[a_0\ot \ldots\ot a_n]$ of $[0]$'s and $[1]$'s (resp.\ and $\emptyset$) to the sum of the two sequences that result from adding a $[0]$ or a $[1]$ in position $i$ and applying $\rho^{-1}$ to the terms in position $\geq i$. The coface map $\d^i$ of $\Per{r}$ sends $\sus{q}[0]$ to $\sus{q+1}[0]$ and $\sus{q}[1]$ to $\sus{q+1}[1]$. Thus $\psi$ commutes with $\d^i$, and therefore is a cosemi-simplicial chain map.
\begin{remark}
	The article \cite{medina2021fast_sq} predates the viewpoint of this one in the following way: The maps $e_q^\dd\ot x \mapsto \nabla^{\binom{n}{q}}(x\ot x)$ define a cyclic connected homotopy comultiplication on the chain complex of a simplicial set. 
\end{remark}

\begin{remark}	
	The combinatorics in pages 13 and 14 \federico{check pages, I don't have access to the published version} of \cite{medina2021fast_sq} translate here to the fact that the map from the unnormalized chains to the normalized chains is well-defined, i.e., to the fact that the degenerate simplices generate a subcomplex of the unnormalized chain complex. To add more detail: Medina-Mardones needs to check that the map $\alpha^\vee\circ \beta^{\vee}\circ \psi^\vee\circ \phi^\vee$ (denoted $\nabla^{\binom{n}{q}}$) is a cosemi-simplicial chain map, though most of the trouble is just proving that it is a chain map.
	
	As we have seen in the previous subsection, the map $\psi$ factors as $\mathrm{red}\circ \psi_{\mathrm{red}}$ and we have
	\[
		e_q^{\vee}\overset{\phi^\vee}{\longmapsto}
		\susp{q}(0)\overset{\psi^\vee_{\mathrm{red}}}{\longmapsto}
		(0)\ot(1)\ot(0)\overset{q}{\ldots} 
	\]
	Now, the map $\mathrm{red}$ takes this last generator and inserts $n+1-q$ instances of $\emptyset$ in all possible ways, for example, if $q=2$ and $n=3$, then we obtain the following summands:
	\begin{align*}
		&(0)\ot(1)\ot\emptyset\ot\emptyset
		&
		&(0)\ot\emptyset\ot(1)\ot\emptyset
		&
		&(0)\ot\emptyset\ot\emptyset\ot(1)
		\\
		&\emptyset\ot(0)\ot(1)\ot\emptyset
		&
		&\emptyset\ot(0)\ot\emptyset\ot(1)
		&
		&\emptyset\ot\emptyset\ot(0)\ot(1)
	\end{align*}
	Each of these summands can be indexed by the positions of the non-empty entries
	\begin{align*}
		&(0,1)&& (0,2)&&(0,3)
		\\
		&(1,2)&& (1,3) &&(2,3)
	\end{align*}
	These are the elements $U\in \cP_q(n)$ in Medina-Mardones' notation. The map $\beta^\vee$ sends these summands to
	\begin{align*}
		&(0)\ot(0)\ot\emptyset\ot\emptyset
		&
		&(0)\ot\emptyset\ot(1)\ot\emptyset
		&
		&(0)\ot\emptyset\ot\emptyset\ot(0)
		\\
		&\emptyset\ot(1)\ot(1)\ot\emptyset
		&
		&\emptyset\ot(1)\ot\emptyset\ot(0)
		&
		&\emptyset\ot\emptyset\ot(0)\ot(0)
	\end{align*}
	and the map $\alpha$ sends these summands to
	\begin{align*}
		&(0,1)\ot \emptyset
		&
		&(0)\ot (2)
		&
		&(0,3)\ot \emptyset
		\\
		&\emptyset\ot(1,2)
		&
		&(3)\ot (1)
		&
		&(2,3)\ot\emptyset
	\end{align*}
	which correspond to the decomposition $U^0\ot U^1$ of each $U$ in Medina-Mardones' notation. 

	The cancellations taking place in the forementioned pages are due to the fact that $\psi_{\mathrm{red}}^\vee \circ \phi^\vee(e_q^\vee)$ in $\cochains(\partial\simplex^{r-1})^{\ot n+1}\cong \ucochains(\EC_2)$ is a normalised cochain, thus its coboundary is again a normalised cochain. For example, in $\cochains(\asimplex^{n})^{\ot r}_{\interior}$ the coboundary of $(0)\ot (2)$ has a summand $(0,1)\ot (2)$ that also appears in the coboundary of $(0,1)\ot\emptyset$, thus they cancel out. If we translate this back to the chain complex of the reduced join $\cochains(\partial\simplex^{r-1})^{\ot 4}$, we see that $(0)\ot (2)$ and $(0,1)\ot\emptyset$ are summands in the image of $(0)\ot (1)$ and that $(0,1)\ot (2)$ is a summand in the image of $(0)\ot (1)\ot(1)$. Now, in the coboundary of $(0)\ot (1)$ the summand $(0)\ot (1)\ot(1)$ appears twice: one comes from inserting a $(1)$ in the middle position and another from inserting a $(1)$ in the right position, thus they cancel out.
\end{remark}

\subsection{Kan spectra} \label{s:9Kanspectra}

The category of Kan spectra is one of the many categories that model spectra. Introduced by Kan in \cite{Kan1963} and later developed by several articles \cite{burghelea_kanspectraI1967, burghelea_kanspectraII_1968, burghelea_kanspectraIII1969,Brown1973} (see alsi \cite{Stephan2015, CKP2023}).

\begin{definition}
	The \emph{stable simplex category} $\kansimplex$ is the colimit of the functors
	\[
	\simplex \lra \simplex \lra \simplex \lra \dots
	\]
	that send an ordinal $[n]$ to the ordinal $[n+1]$ and extend a map $f \colon [n] \to [m]$ to a map $f' \colon [n+1] \to [m+1]$ by setting $f'(n+1) = m+1$. It has one object $[n]$ for each integer $n\in \bZ$, and there are morphisms $\d^i \colon [n-1] \to [n]$ and $\s^i \colon [n] \to [n-1]$, for all $i \geq 0$, satisfying the simplicial identities.
\end{definition}

The morphisms from $[n]$ to $[m]$ are in bijection with the endomorphisms $f$ of the first infinite ordinal $\omega$ such that
\begin{itemize}
	\item $|f^{-1}(k)|<\infty$ for all $k\in \omega$,
	\item $|f^{-1}(k)|\neq 1$ for finitely many $k\in \omega$,
	\item $\sum_k (1-f^{-1}(k)) = m-n$.
\end{itemize}

\begin{definition}
	A \emph{Kan spectrum} is a presheaf $X \colon \kansimplex^{\op} \to \Setp$ on pointed sets such that for each $x\in X_n$ there is an $m$ such that $d_i(x) = *$ for all $i>m$. A map of Kan spectra is a natural transformation.
\end{definition}

The maps $\simplex\to \simplex$ induce chain maps $\chains(\asimplex)\to \sus{}\chains(\asimplex)$ and after taking Alexander duality, chain maps $\cochains(\asimplex)\to \cochains(\asimplex)$. The cochain complex of the standart $\infty$-simplex is $\cochains(\asimplex^\infty)\cong \bigcup_n \cochains(\asimplex^n)$.


The isomorphism \eqref{eq:1} composed with \eqref{eq:dual} yields an isomorphism
\[
		\susp{\bullet}\cochains(\asimplex^{\infty}) \ot_{\kansimplex} X_\bullet \cong \chains(X).	
\]
The rest of the construction applies verbatim replacing
\[
\susp{\bullet}\cochains(\asimplex^{\bullet})
\quad \text{by}
\quad  \susp{\bullet}\cochains(\asimplex^{\infty})
,\] 
and
\[
\sus{\bullet(r-1)}\chains(\asimplex^r)^{\ot \bullet}
\quad \text{by} \quad
\sus{\bullet(r-1)}\bigcup_n\chains(\asimplex^n)^{\ot n}.
\]
%
%
%This condition allows to define chain complexes $\ucadenas(X)$ and $\cadenas(X)$ as for simplicial sets: In degree $n$ they are generated by the pointed set $X_n$ modulo the base point, and the differential is the alternate sum of the face maps.
%
%The standard augmented infinite simplex $\asimplex$ is the Kan spectrum with one $n$-simplex for each order-preserving map $[n] \to [\omega]$.
%
%Write as before $\uchains(X)$ and $\chains(X)$ for their shifted versions. There is a covariant functor $\kansimplex \to \Ch{R}$ that sends any object $[n]$ to the chain complex of the augmented infinite simplex $\chains(\asimplex^\infty)$, and the action of faces and degeneracies are the stabilization of the faces and degeneracies of the finite simplices. There is another covariant functor $\kansimplex \to \Ch{R}$ that sends any object $[n]$ to the cochains $\cochains(\asimplex^\infty)$, and the action of the faces and degeneracies is the limit of their finite counterparts (note that this second complex is not free). Then, we have, as in Lemma \ref{lemma:2}:
%
%\begin{lemma}
%	There is a natural isomorphism
%	\[
%	\chains(\asimplex^{\infty}) \ot_{\kansimplex} \rA(X)_\bullet \cong \cadenas(X)
%	\]
%\end{lemma}
%
%Then, the rest of \cref{s:3complexes} remains true: One obtains a stable cosimplicial cochain complex $\chains(\asimplex^\infty)$ and a stable cosimplicial cochain complexes $\Omd(r,\infty)$ and $\Omhatd(r,\infty)$, and the former is the dual to a chain complex $\Om(r,\infty)$ that admits maps $\eta^{n}$ to $\Om(r)$. With these changes, Proposition \ref{prop:omegarm} yields a connected diagonal on the chain complex of a Kan spectrum.
%
%Note that the previous lemma does not make sense without dualizing: the top simplex in $N_*(\asimplex^\infty)$ is far away, at infinite dimension.

\subsection{Behaviour of the May--Steenrod structure on simplicial sets under suspension}\label{s:suspension} 
\begin{proposition}\label{prop:suspensionunstable}
	The effect of left and right suspensions on the $r$-cyclic unstable diagonals of \cite{medina2021may_st} is the following:
	\begin{align*}
		\mu(e_q \ot \susp{} \tau) &= (-1)^{q+\binom{r}{2}n}(\tilde{r}-1)!\susp{\ot r}\sum_{i=1}^{\tilde{r}} \rho^{2i}\mu(e_{q-r+1} \ot \tau)
		\\
		\mu(e_q \ot \sus{}\tau) &= (-1)^{q+\binom{r}{2}n}\tilde{r}!\sus{\ot r}\mu(e_{q-r+1} \ot \tau)
	\end{align*}
\end{proposition}

\begin{proof}
	The effect of suspension on the Barrat-Eccles operad and the surjection operad is described in the appendix of \cite{berger2004combinatorial}. From that description one deduces that the summands on the left corresponding to an element $(\rho^0,\rho^{i_1},\rho^{i_1+1},\dots,\rho^{i_k},\rho^{i_k+1})$ will vanish unless the first $r$ elements are all distinct. A count yields the result.
	
	For the second statement, we follow the same argument, except that the summands that do not vanish are those whose last $r$ elements are all distinct.
\end{proof}

%If $\mu$ is the unstable diagonal of \cite{medina2021may_st}, there is a stable diagonal $\mu_{st}$ on the stable chains of a $\sus{}$-spectrum with coefficients in $\bZ[\frac{1}{\tilde{r}!}]$ defined as
%\[
%\mu_{st}(e_q \ot \tau) = (-1)^{mq+\binom{r}{2}\left(m(m+n-1) + \binom{m+1}{2}\right)}\frac{1}{(\tilde{r}!)^{m}}\mu(e_{q+m(r-1)}) \ot \tau)
%\]
%where $\tau\in \chains[n](E_m)$. Since $\sum_{i=1}^{\tilde{r}} \rho^{2i}$ is invertible, we have can derive a formula for the left suspension too.

The disparity in the behaviour with respect to left suspension of the operations in this work and the ones in \cite{medina2021may_st} show that

\begin{corollary}
	The chain level operations of \cite{medina2021may_st} are not the same as the ones of this work if $r\geq 5$.
\end{corollary}




\subsection{Effective computations} \label{s:python} The maps in this work have been implemented in the repository \cite{github-oddp}. 

\subsubsection{The abstract comultiplication} The following are the times (in seconds) required to find $\alpha^{\dd}\circ \beta^{\dd}\circ \mathrm{red}\circ \psi_{\mathrm{red}}^\dd\circ \phi^\dd(e_{q}^{\dd})$ on dimension $n$ for $r=3$. Notice that the maps $\alpha$ and $\beta$ are isomorphisms. We have observed that most of the time is used to compute the reduction map $\mathrm{red}$:

For $q=3$ (this computes the Bockstein):
\[
\begin{array}{|c|c|c|c|c|c|c|c|c|c|c|} \hline
	n & 10  & 20    & 30    & 40   & 50   & 60 & 70 & 80 & 90 & 100
	\\ \hline
	  &0    & 0.04  & 0.11  & 0.32 & 0.71 & 1.5 & 2.5& 4.2 & 6.7 & 8.3 \\ \hline
\end{array}
\]For $q=12$ (this computes the first non-trivial operation after the Bockstein):
\begin{equation}\label{eq:comp}
\begin{array}{|c|c|c|c|c|c|c|c|c|c|c|c|} \hline
	n & 7  & 8    & 9    & 10   & 11   & 12 & 13 & 14 & 15 & 16 & 17
	\\ \hline
	&0 &0.02& 0.02 & 0.07 & 0.27 & 0.85 & 2.7 & 7.5 & 20 & 52 & 121 \\ \hline
\end{array}
\end{equation}
For $q= 15$ (this computes the Bockstein of the previous operation)
\[
\begin{array}{|c|c|c|c|c|c|c|c|c|c|c|c|} \hline
	n & 8    & 9    & 10   & 11   & 12 & 13 & 14 & 15 & 16 & 17
	\\ \hline
	&0 & 0      & 0.03 & 0.09 & 0.4 & 1.7 & 6 & 21 & 65 & 204 \\ \hline
\end{array}
\]

\subsubsection{Comparison with comch} Computing the formulas of the table \eqref{eq:comp} with comch takes the following times\federico{The comparison is a bit unfair, because comch might be optimized to give better times when computing power operations}:
\[
\begin{array}{|c|c|c|c|} \hline
	n & 7 & 8 & 9  \\ \hline
	&0  & 0.5 & 38 \\ \hline
\end{array}
\]

\subsubsection{Evaluating the formulas on particular simplicial complexes}
Below are the times taken to compute the Bockstein ($q=3$) for the model of the Moore space $M(\bZ/3)$ of sagemath and its $m$-fold unreduced suspension (that adds two vertices and multiplies the number of faces by $3^m$). Below we write $n=m+3$ for the degree of the top simplex in $\susp{m}M(\bZ/3)$, to ease comparison with the previous tables.
\[
\begin{array}{|c|c|c|c|c|c|c|c|c|c|c|c|c|c|c|}\hline 
	n &3 & 4 & 5 & 6 & 7 & 8 & 9  & 10 & 11    & 12 & 13 & 14 & 15\\ \hline
	  &0 & 0  & 0 & 0 & 0 & 0 & 0.03  & 0.1 & 0.25 & 0.7 & 2 & 5 & 15\\ \hline
\end{array}
\]
%Curiously the next column gives memory error.
We then implemented a small simplicial model of $\bC P^3$ found by \cite{datta} with 18 vertices, and considered its $m$th unreduced suspension (that adds two vertices and multiplies the number of faces by $3^m$) and applied to it an algorithm to compute the first non-trivial operation after the Bockstein ($q=12$). Below we write $n=m+7$ for the degree of the top simplex in $\susp{m}\bC P^3$, to ease comparison with the previous tables.
\[
\begin{array}{|c|c|c|c|c|}\hline 
	n & 7     & 8    & 9 & 10 \\ \hline
	  & 0.02  & 0.48 & 7 & 87 \\ \hline
\end{array}
\]
\federico{One can also take the reduced suspension, in which case the number of faces stays constant and the performance is better. Maybe we can write both.}


The computer used in these calculations had a 13th Gen Intel(R) Core(TM) i5-1335 1.30GHz and 8GB of RAM. %(13th Gen Intel® Core™ i5-13420H × 12 and 16GB RAM)